Gom các bài lập trình nộp sai theo những test case mà chúng trượt, rồi giải thích mỗi cụm bằng luật NẾU--THÌ, là cách phổ biến để cho giảng viên biết cả lớp đang có chung những quan niệm sai nào. Việc các cụm đó có gom chương trình theo cơ chế lỗi hay không hiếm khi được kiểm chứng, vì thiếu nhãn cơ chế. Chúng tôi xây các nhãn này mà không cần gán nhãn thủ công. Sau khi chạy lại {{n_submissions}} bài nộp C90 của một corpus công khai trong runner cách ly (khớp {{audit_cells_pct}}\% với hệ chấm lịch sử), chúng tôi thu gọn mỗi bài sai và lần nộp đạt kế tiếp của cùng sinh viên thành một bản sửa tối thiểu đã kiểm chứng bằng delta debugging, rồi xếp nó vào một trong mười hai cơ chế mức chương trình ({{n_real_gold}} nhãn đơn cơ chế); một benchmark thứ hai tiêm đúng một cơ chế vào các chương trình đã đạt ({{n_injected}} mutant). Theo protocol đăng ký trước, các chữ ký test giống hệt trộn lẫn cơ chế: tính trung bình trên các cohort bài tập, không hàm nào của biểu diễn kết quả test xếp đúng được quá {{real_H1a_pct}}\% item thật và {{injected_H1a_pct}}\% item tiêm lỗi. Các mô tả không cần nhãn về cách output lệch khỏi đáp án nâng trần này vượt một mô hình ngẫu nhiên cùng độ mịn và, trên dữ liệu có kiểm soát, nâng chỉ số Rand hiệu chỉnh của cụm K-means từ {{injected_ari_outcomes}} lên {{injected_ari_deviation}}; trên bản sửa thật mức tăng ({{real_ari_outcomes}} lên {{real_ari_deviation}}) không đáng tin cậy, và các mô tả này không vượt được các quan hệ output đơn giản hơn. Luật ILA-2 trên các mô tả đó chuyển được sang bài tập chưa thấy (macro-F1 {{injected_unseen_problems_evidence_ila2_macro_f1}}), trong khi một embedding mã đóng băng gom mutant theo chương trình gốc chứ không theo cơ chế. Một công cụ cho giảng viên dựng trên các kết quả này chỉ nêu giả thuyết cơ chế khi một luật khớp phần lớn một nhóm. Với phản hồi ở mức cơ chế, cải tiến rẻ nhất là bằng chứng mà hệ chấm tự động vốn đã có.
